\subsection{The final case: inverse \texorpdfstring{$V_4$}{V4} Morita duality}
\label{subsec:inverse-v4}
The surviving quadratic $A_4$ ring at $m=8$ cannot reconstruct the
rank-seven source. The obstruction is independent of uniqueness or
unitarity of its categorifications.

\begin{lemma}\label{lem:v4-commutation}
Let $\mathcal E$ be a complex fusion category with invertibles
$H=V_4$, trivial pointed associator, and four other simples $h\rho$,
where
\[
 \rho h=h\rho,\qquad \rho^2=1+\sum_{h\in H}h\rho.
\]
The commutation cocycle of $\rho$ with $H$ is trivial in
$H^2(H,\C^\times)$.
\end{lemma}
\begin{proof}
Choose coherent multiplication maps for $H$ and isomorphisms
$h\rho\to\rho h$. Their scalar defect is a $2$-cocycle $\xi$;
its cohomology class is independent of these choices
\cite[Section~2]{Izumi3n}. Let $P=\bigoplus_{h\in H}h$ be the
regular algebra. In the bicategory of separable algebras, write
$\iota$ for the canonical module and $\bar\iota$ for its adjoint,
so $\bar\iota\iota=\bigoplus_hh$, and put
$Z=\iota\rho\bar\iota$. Moving the $H$ labels past $\rho$
identifies
\[
              \operatorname{End}(Z)\cong\C_\xi[H].
\]
This is the algebraic calculation in
\cite[Proposition~9.6]{Izumi3n}; it requires only associativity and
rigidity, not a $C^*$ structure.

If $[\xi]\ne0$, the twisted group algebra is $M_2(\C)$, hence
$Z=2\sigma$ for a simple $\sigma$. The pointed Morita dual is
$\operatorname{Rep}(H)=\mathrm{Vec}_{\widehat H}$. The noninvertible
orbit has just this one projective stabilizer representation, and
\[
 Z^2=4\iota\rho^2\bar\iota
     =4\sum_{\chi\in\widehat H}\chi+16Z.
\]
Consequently
$\sigma^2=\sum_{\chi\in\widehat H}\chi+8\sigma$.
This would categorify the near-group ring $R(V_4,8)$, whereas
\cite[Theorem~1.1]{Schopieray} allows only levels $0$ and $4$
for $V_4$. That theorem concerns arbitrary complex categorification.
Replacing the $C^*$-only final citation of Izumi's argument by this
general result proves the lemma in the stated scope.
\end{proof}

\begin{theorem}\label{thm:r1022-general}
The ring $R7\_3\_1022$ has no categorification over $\C$.
\end{theorem}
\begin{proof}
Its dimensions are $1,1,1,3,d,d,d$, where
$d=4+2\sqrt5$ and $d^2=4+8d$. The two Morita transformations of
Theorem~\ref{thm:c3-a4}, with $m=8$, give
\[
 \mathcal D:\quad 1^{12},d^3,
 \qquad
 \mathcal D'=\mathcal D_B^*:\quad 1^{12},t^{12},
 \qquad t=d/2=2+\sqrt5,
\]
where $B$ is one of the regular $V_4$ algebras and both pointed
categories are untwisted $\mathrm{Vec}_{A_4}$. The same quadratic
ring normalization has no outer branch at $k=4$; its inner branches
are $N=4\cdot1$ and $N=\sum_{h\in V_4}h$. The first contains a
forbidden rank-two subcategory. Thus in $\mathcal D'$
\[
 \rho g=g\rho\quad(g\in A_4),\qquad
 \rho^2=1+\sum_{h\in V_4}h\rho.
\]
The subcategory generated by $V_4$ and $\rho$ satisfies
Lemma~\ref{lem:v4-commutation}.

The inverse algebra is forced by the actual Morita context.
In $\mathcal M=\mathcal D_B$, the regular right module $m=B$ is
simple since $\operatorname{End}_B(B)=\operatorname{Hom}(1,B)=\C$.
Let $Q=\underline{\operatorname{End}}_{\mathcal D'}(m)$.
Double centralizer Morita duality recovers $\mathcal D$ as
$(\mathcal D')_Q^*$, up to the harmless opposite convention. For
any simple $B$-bimodule $X$, its multiplicity, or that of $X^*$
according to convention, in $Q$ is
\[
 \dim\operatorname{Hom}_{\mathcal M}(B,B\otimes_BX)
 =\dim\operatorname{Hom}_B(B,X)
 =\dim\operatorname{Hom}_{\mathcal D}(1,X).
\]
Exactly the four invertible bimodules supported on the original
$V_4$ contain $1$. Each has underlying object $B$, with multiplicity
one on $1$, for either cocycle defining $B$. Every other bimodule
has support outside that pointed orbit. Hence $Q$ is the sum of the
four invertibles of the unique subgroup $H'=V_4\subset A_4$,
with one of its two algebra cocycles, say $\beta$.

The twelve noninvertibles of $\mathcal D'$ divide into three free
$H'$-orbits of size four. At $g\rho$ the left-right stabilizer is
the graph of conjugation by $g$ on $H'$. Its projective class is
the transported commutation class, plus
$[\beta]-(\operatorname{Ad}_g)^*[\beta]$.
The first term vanishes by the lemma. The second vanishes because
$H^2(V_4,\C^\times)=C_2$, on which every automorphism acts trivially.
Each orbit therefore has four one-dimensional stabilizer
representations. They give four simple $Q$-bimodules with underlying
dimension $4t$ and dual dimension $t$.

Thus $(\mathcal D')_Q^*$ has twelve noninvertibles of dimension $t$.
It cannot be equivalent to $\mathcal D$, which has three
noninvertibles of dimension $2t$. This contradiction proves the theorem.
\end{proof}
